\section{Collected Works page 551},
%WHEN \cite{rayleigh1902scientific, rayleigh1881}

With bibtex related,
does this work\cite{Rayleigh1881,ScientificPapersVol1,RayleighSoundVol1Ed2}
%or this\cite{Rayleigh1881,RayleighSoundVol1Ed2}
%or is it\cite{ScientificPapersVol1,RayleighSoundVol1Ed2}

When a thin sheet of matter is subjected to stress, the force which it opposes
to extension is great in comparison with that which it apposes to bending.
Under ordinary circumstances, the deformation takes place approximately as if
the sheet were inextensible as a whole, a condition which in a remarkable
degree facilitates calculation, though (it need scarcely be said even bending
implies an extension of all but the central layers.
The inextensibility postulated refers properly to the central layer, which is
spoken of as a surface.
We will commence with the case of a surface naturally spherical, and
investigate what kind of deformation it admits of, under the condition that no
line traced upon the surface is altered in length. The radius of the sphere
being \(a\), let the point whose natural coordinates are \(a\), \(\theta\), \(\phi\)
be displaced to the position
\[ (a + \delta r, \theta + \delta\theta, \phi + \delta \phi),\]
 where \((\delta r,\delta \theta,\delta \phi)\) is to be treated as small.
\[
M(a,\theta) =
\begin{bmatrix}
a^2 & 0
0  & a^2 \sin^2(\theta)
\end{bmatrix}
\]
Since the element of arc \(ds\) is of the same length after displacement as
\[
\norm{ds}^2 = (d\theta, d \phi) M [d\theta; d \phi ]
\]
I am just making up some notation here;  why not \(x\) ?
\[
x^T = (d\theta + d \delta \theta, \phi + d \delta \phi ),
\]
Inextensibility means
\[
\norm{ds}^2 = x^T  M(a +  \delta r,\theta + \delta\theta) x
\]
Rayleigh includes the addition term \( (d \delta r)^2 \); I do not see
where this comes from; I will just assume that it is \(0\), and see
how far I get.


What is \(x\)?
\[
\Grad = (\frac{\del}{\del \theta}, \frac{\del}{\del \phi}),
\]
\[
df = (\Grad f)  [d \theta; d \phi],
\]
\[
x^T = (d\theta + d \delta \theta, d \phi + d \delta \phi ),
\]
\begin{align}
d\theta + d \delta \theta = ((1,0) +\Grad (\delta \theta)) [d \theta; d \phi] \\
d \phi + d \delta \phi  = ((0,1) + \Grad (\delta \phi)) [d \theta; d \phi]
\end{align}
\[
x = \left( I +
D \begin{matrix}
d\theta &
d \phi
\end{matrix} \right)
\begin{matrix}
d\theta &
d \phi
\end{matrix}
\]
\[
= (\frac{\del f}{\del \theta}, \frac{\del f}{\del \phi}
(\frac{\del f}{\del \theta}, \frac{\del f}{\del \phi}
\]
\[ d f = \frac{\del f}{\del \theta} d \theta  + d f = \frac{\del f}{\del \theta} d \theta \]
\[ d \delta \theta =(\Grad \delta \theta )[d \theta;d \phi]
\frac{ \del \delta \theta }{ \del \theta } d \theta
\]

\begin{verbatim}
= (d\theta, d \phi) M(a +  \delta r,\theta + \delta\theta) [d\theta; d \phi ]
, \phi + d \delta \phi ).
 + d \delta ):,
 + d \delta \theta) + (d \delta r) + (a + \delta r) sin' (\theta + \delta\theta)
(d \phi + d \delta ):,

= (a + \delta r)
(d\theta + d \delta \theta) + (d \delta r) +
(a + \delta r) sin' (\theta + \delta\theta)
(\phi + d \delta \phi ):,
\end{verbatim}
before, we have


This is raw text from an old ocr scanner.
\begin{verbatim}
= (a + \delta r) (d\theta + d \delta \theta) +
(d \delta r) +
(a + \delta r) sin' (\theta + \delta\theta) (2\phi + d \delta ):,

or, retaining only the first power of \delta r, \delta\theta, \delta\phi,


 + ir (de) + ir sin e (d theta) +
a sin \theta cos \theta \delta\theta (d)* +
a sin* \theta dd = \theta.
\end{verbatim}
Now
\[
= (a + \delta r) (d\theta +
d \delta \theta) + (d \delta r) +
(a + \delta r) sin' (\theta + \delta\theta) (2\phi + d \delta ):,
\]
just write e for delta ?
